\section{Introduction}

Migration has substantially changed the composition of the population in Germany over the past decade. Two groups are particularly important in this development: migrants from Syria and Ukraine. Syrian migration to Germany increased substantially during the refugee inflow of 2015 and 2016. In 2015 alone, approximately 326,900 Syrian nationals arrived in Germany \citep{bamf2015migration}. More recently, Russia's full-scale invasion of Ukraine in February 2022 led to another large movement of people seeking protection in Germany. Following the activation of the European Union's Temporary Protection Directive, people fleeing Ukraine could obtain temporary protection in Germany without going through the regular asylum procedure. The number of Ukrainian nationals living in Germany consequently increased from approximately 155,000 at the end of 2021 to more than 1.16 million at the end of 2022 \citep{bamf2022migration}. By the end of 2025, around 1.41 million Ukrainian and 936,000 Syrian nationals were registered in Germany, making them two of the country's largest foreign-national populations \citep{destatis2026foreign}.

The prominence of these two groups makes Germany a particularly relevant setting for examining how beliefs about migrants are formed. Ukrainian and Syrian migrants arrived in Germany in different historical and geopolitical contexts and through partly different migration pathways. At the same time, both groups include individuals with heterogeneous characteristics, including differences in education, language proficiency, reason for migration, and labor-market experience. Comparing beliefs about Ukrainians and Syrians therefore raises an important question: to what extent do differences in beliefs reflect country of origin itself, and to what extent are they related to other characteristics of migrants?

Previous research shows that attitudes toward immigrants and refugees depend on characteristics such as country of origin, education, economic prospects, and reason for migration \citep{bansak2016economic, czymara2016who}. The arrival of Ukrainian refugees has further renewed interest in differences in the reception and perception of refugee groups. Recent research documents differences in European attitudes toward Ukrainian and other refugees and emphasizes the importance of the broader political and migration context in which these attitudes are formed \citep{moise2024european}. Research comparing the political framing of Syrian and Ukrainian refugees likewise demonstrates that the two refugee movements occurred in distinct contexts, while Germany represents a particularly informative case because it received substantial numbers of refugees from both groups \citep{drewski2025states}.

However, observed differences in beliefs about migrant groups are difficult to interpret when the groups also differ along other dimensions. A Syrian and a Ukrainian migrant encountered in everyday life may differ not only in country of origin, but also in education, German-language proficiency, migration reason, gender, or other characteristics. Moreover, beliefs about migrants may not depend only on the characteristics of the migrants themselves. They may also be related to respondents' previous contact and experiences with different migrant groups, the information they encounter through the media, and the experiences or information that remain particularly memorable. Understanding differences in beliefs therefore requires considering both the characteristics of the migrant being evaluated and the information and experiences available to the person forming the belief.

Against this background, this thesis examines how individuals in Germany form beliefs about the labor-market and social integration of Ukrainian and Syrian migrants. It asks whether otherwise comparable Ukrainian and Syrian migrants are expected to achieve different integration outcomes and how these beliefs depend on other migrant characteristics, particularly human capital, German-language proficiency, and reason for migration. The thesis further examines how respondents' beliefs are associated with their personal contact and experiences with Ukrainian and Syrian migrants, their exposure to migration-related media information, and the valence of memorable personal experiences and media information.

To address these questions, the thesis develops a factorial survey experiment in which respondents evaluate hypothetical migrant profiles. Country of origin, reason for migration, human capital, German-language proficiency, and gender are independently randomized across profiles. Respondents report their beliefs about each migrant's employment prospects, future income, exposure to hiring and workplace discrimination, and social integration. This design allows the causal effect of experimentally varied migrant characteristics on stated beliefs to be identified while holding the other presented characteristics constant. The vignette experiment is complemented by a respondent questionnaire measuring personal contact and experience, media exposure and tone, memorable experiences and information, and broader beliefs about migrant communities and integration over time. Since these respondent-level measures are not experimentally assigned, the corresponding analyses are interpreted as associations rather than causal effects.

The contribution of the thesis is therefore twofold. First, it provides a framework for separating country-of-origin differences in beliefs from differences associated with other migrant characteristics. This is particularly relevant for the comparison of Ukrainian and Syrian migrants, for whom country of origin, refugee status, human capital, and other characteristics may otherwise be difficult to disentangle. Second, the thesis connects the experimental analysis of migrant characteristics with respondents' personal experiences and information environments. In doing so, it brings together research on preferences and beliefs about migrant characteristics with research emphasizing contact, information exposure, and memory in belief formation.

The empirical study developed in this thesis is proposed rather than implemented. The survey instrument, randomization procedure, variable construction, and empirical specifications are therefore described in sufficient detail to provide an implementation-ready design that can be conducted in future research. Accordingly, the thesis does not report empirical results but develops a pre-specified framework for testing the proposed hypotheses and for distinguishing causal experimental effects from associational evidence on respondents' experiences and information exposure.

The remainder of the thesis is structured as follows. Section 2 reviews the relevant literature and develops the hypotheses. Section 3 presents the proposed experimental design, outcome measures, respondent questionnaire, and data-processing procedures. Section 4 describes the empirical strategy for testing the experimental and associational hypotheses and the additional exploratory analyses. Section 5 discusses the main limitations of the proposed design and directions for future research. Section 6 concludes.
